\subsection{AN EXAMPLE}

Let \(k\) be a field, \(n\) an integer \(\geqslant 0\) , and \((e_i)\) the canonical basis of \(k^n\) . Let \(G = GL(n, k)\) , let \(B\) be the upper triangular subgroup of \(G\) , and let \(N\) be the subgroup of \(G\) consisting of the matrices having exactly one non-zero element in each row and column. An element of \(N\) permutes the lines \(ke_i\) ; this gives rise to a surjective homomorphism \(N \to \mathfrak{S}_n\) whose kernel is the subgroup \(T = B \cap N\) of diagonal matrices, and allows us to identify \(W = N/T\) with \(\mathfrak{S}_n\) . We denote by \(s_j\) ( \(1 \leqslant j \leqslant n - 1\) ) the element of \(W\) corresponding to the transposition of \(j\) and \(j + 1\) ; let \(S\) be the set of \(s_j\) . The quadruple \((G, B, N, S)\) is a Tits system. Indeed:

Axiom (T1) follows from Cor. 2 of Prop. 14 of Algebra, Chap. II, § 10, no. 13.

Axiom (T2) is proved in Algebra, Chap. I, Correction to p.~97.

Axiom (T4) is immediate.

It remains to verify axiom (T3), i.e.

\[
s _ {j} \mathrm{B} w \subset \mathrm{B} w \mathrm{B} \cup \mathrm{B} s _ {j} w \mathrm{B} \quad \text { for } 1 \leqslant j \leqslant n - 1, w \in \mathrm{W},
\]

or equivalently,

\[
s _ {j} \mathrm{B} \subset \mathrm{BB} ^ {\prime} \cup \mathrm{Bs} _ {j} \mathrm{B} ^ {\prime}, \quad \text { with } \mathrm{B} ^ {\prime} = w \mathrm{B} w ^ {- 1}.
\]

Let \(G_{j}\) be the subgroup of G consisting of the elements that fix the \(e_{i}\) for \(i \neq j, j + 1\) and stabilize the plane spanned by \(e_{j}\) and \(e_{j+1}\) ; this group is isomorphic to \(\mathbf{GL}(2, k)\) . One checks that \(G_{j}B = BG_{j}\) . Since \(s_{j} \in G_{j}\) , we have \(s_{j}B \subset BG_{j}\) , and it suffices to prove that

\[
\mathrm{G} _ {j} \subset (\mathrm{B} \cap \mathrm{G} _ {j}) (\mathrm{B} ^ {\prime} \cap \mathrm{G} _ {j}) \cup (\mathrm{B} \cap \mathrm{G} _ {j}) s _ {j} (\mathrm{B} ^ {\prime} \cap \mathrm{G} _ {j}).
\]

Identify \(G_{j}\) with \(\mathbf{GL}(2,k)\) ; the group \(B \cap G_{j}\) is then identified with the upper triangular subgroup \(B_{2}\) of \(\mathbf{GL}(2,k)\) , while the group \(B' \cap G_{j}\) is identified with \(\mathbf{B}_2\) when \(w(j) < w(j + 1)\) and with the lower triangular subgroup \(\mathbf{B}_2^-\) otherwise. In the first case, the formula to be proved can be written

\[
\mathbf {G L} (2, k) = \mathrm{B} _ {2} \cup \mathrm{B} _ {2} s \mathrm{B} _ {2} \quad \text { where } \quad s = \left( \begin{array}{c c} 0 & 1 \\ 1 & 0 \end{array} \right);
\]

this follows for example from the fact that \(B_{2}\) is the stabilizer of a point for the action of \(\mathbf{GL}(2,k)\) on the projective line \(\mathbf{P}_{1}(k)\) , and acts transitively on the complement of this point. In the second case, the formula to be proved can be written

\[
\mathbf {G L} (2, k) = \mathrm{B} _ {2} \mathrm{B} _ {2} ^ {-} \cup \mathrm{B} _ {2} s \mathrm{B} _ {2} ^ {-};
\]

since \(B_{2}^{-} = sB_{2}s\) , this follows from the preceding formula by multiplying on the right by s.

